% Generated draft: requires mathematical and source review.
\documentclass[11pt]{article}
\usepackage[margin=25mm]{geometry}
\usepackage{fontspec}
\setmainfont{Latin Modern Roman}
\usepackage{amsmath,amssymb,mathtools,unicode-math}
\setmathfont{Latin Modern Math}
\usepackage{enumitem,needspace,xurl,hyperref}
\setlength{\parindent}{0pt}
\setlength{\parskip}{5pt}
\setlength{\emergencystretch}{3em}
\newcommand{\N}{\mathbb{N}}
\newcommand{\Z}{\mathbb{Z}}
\newcommand{\Q}{\mathbb{Q}}
\newcommand{\R}{\mathbb{R}}
\newcommand{\C}{\mathbb{C}}
\newcommand{\F}{\mathbb{F}}
\newcommand{\A}{\mathbb{A}}
\newcommand{\PP}{\mathbb{P}}
\newcommand{\E}{\mathbb{E}}
\newcommand{\eps}{\varepsilon}
\newcommand{\dd}{\,\mathrm{d}}
\begin{document}
\section*{UnsolvedMath ID: 800708}
\textbf{Source problem code:} catalogue entry
\subsection{Definitions and mathematical statement}
Let \(F\) be a field. For positive integers \(d\mid n\), let \(\mu_d\) be the group scheme of \(d\)-th roots of unity over \(F\), embedded centrally in \(\operatorname{GL}_n\) by scalar matrices, and let \(\operatorname{GL}_n/\mu_d\) denote the quotient algebraic group. The group scheme interpretation applies also when \(\operatorname{char}F\mid d\); replacing \(\mu_d\) by its group of \(F\)-rational points would change the meaning.

For an algebraic group \(G\) over \(F\) and a field extension \(K/F\), let \(H^1_{\mathrm{fppf}}(K,G)\) denote the set of isomorphism classes of \(G_K\)-torsors. For \(\xi\in H^1_{\mathrm{fppf}}(K,G)\), define
\[
\operatorname{ed}_F(\xi)=\min\left\{\operatorname{trdeg}_F L:\ F\subseteq L\subseteq K,\quad \xi\in\operatorname{im}\bigl(H^1_{\mathrm{fppf}}(L,G)\longrightarrow H^1_{\mathrm{fppf}}(K,G)\bigr)\right\}.
\]
The essential dimension of \(G\) over \(F\) is
\[
\operatorname{ed}_F(G)=\sup_{K/F}\ \sup_{\xi\in H^1_{\mathrm{fppf}}(K,G)}\operatorname{ed}_F(\xi).
\]

Determine the exact values of the following essential dimensions:
\begin{enumerate}
\item \(\operatorname{ed}_F(\operatorname{GL}_{8}/\mu_2)\) for every field \(F\) with \(\operatorname{char}F=2\).
\item \(\operatorname{ed}_F(\operatorname{GL}_{16}/\mu_2)\) for every field \(F\) with \(\operatorname{char}F\ne2\).
\item \(\operatorname{ed}_F(\operatorname{GL}_{p^2}/\mu_p)\) for every odd prime \(p\) and every field \(F\), including fields of characteristic \(p\).
\end{enumerate}
No further restrictions on \(F\) are supplied. The registry labels this problem open; that label is supplied source metadata, not an independently verified assessment of its present mathematical status.
\subsection{Short English statement}
Determine the essential dimensions of the three specified families of quotient groups over every field in the stated characteristic ranges. Essential dimension measures the number of algebraically independent parameters needed to define torsors. The quotients use roots of unity as group schemes, including in characteristics dividing their orders.
\subsection{Sources}
\begin{itemize}
\item Local UnsolvedMath registry, permanent ID \(800708\), source code \texttt{catalogue entry}, title \emph{Open Problems on Central Simple Algebras --- Problem 18}. The record attributes the problem to Auel, Brussel, Garibaldi, and Vishne, with proposed year 2010 and supplied status \texttt{open}.
\item Auel, Brussel, Garibaldi, and Vishne, \emph{Open Problems on Central Simple Algebras} (2010), Problem 18. Supplied source URL: \url{https://arxiv.org/abs/1006.3304}. The source record reports access on 2026-10-16 and extraction by \texttt{arxiv\_latex}; these are supplied provenance details, not a claim of independent source access or verification.
\end{itemize}
\end{document}
